\begin{abstract}
\noindent
When many LLM agents share a wiki, notes or memory, one agent's mistake can reach others simply by being copied, and after an incident the operator asks who passed what to whom. We show that ordinary logs usually cannot say: they record what agents wrote, not what they read. In a real shared-wiki incident, even a perfect investigator could name the true source of at most \bestTop{} of copied behaviours, and the obvious detector (``the same string appeared twice'') made \naiveCalls{} copy calls of which only \trustedCalls{} survive base-rate calibration.
We prove ceilings on what any tracer can recover from log content, a bound on false copy calls under monoculture, and soundness of a gateway design in which every served copy carries its own load-bearing token: a \emph{trap street}, after the fictitious streets mapmakers draw to catch copiers. We turn the theory into a traceability audit, the gateway, and a benchmark with a hidden read log.
In \nRuns{} offline agent runs of eight model families, tokens identify the immediate source for \benchTagsTop{} of copying agents against \benchEarliestTop{} for the best simple edit-log rule (false accusations \benchTagsFA{} against \benchRuleFA{}, the latter partly by construction). Whether a token survives depends on whether a model copies links verbatim or re-derives access, so the gain ranges from \liftMin{} to \liftMax{} points. Closing the alternative route raised the traceable share of the \nBypassWord{} models that bypassed tokens from \gatedOpenLo{}--\gatedOpenHi{} to 100\% with no loss of task completion, and tracing a planted bad tip cut the re-check list from \recheckAfter{} to \recheckTrace{} of outputs. We also report what failed: three of our own experiment designs were wrong and one hypothesis we stated in advance was not met.
Code, data and benchmark: \ph{URLs}.
\end{abstract}
